\documentclass[11pt]{article}
\usepackage[a4paper,margin=25mm]{geometry}
\usepackage{amsmath,amssymb,amsthm,mathtools}
\usepackage[hidelinks]{hyperref}
\newtheorem{theorem}{Theorem}
\newcommand{\C}{\mathbb C}
\newcommand{\diag}{\operatorname{diag}}
\newcommand{\ord}{\operatorname{ord}}
\title{Independent mathematical review of the finite original-period zero}
\author{Programme mathematical review}
\date{20 September 2026}
\begin{document}
\maketitle

\section{Exact object and scope}
This review checks the complete coefficient calculation in
\texttt{certify\_finite\_period\_zero.py}, using PCL1--10,
NC1--5 and OFM1--11 of the original supplied programme sources.
It independently proves a coefficient majorant, the local matrix
inverse, the complex zero count and the five-quadric orbit test.
The parameters throughout are
\[
 \delta=\frac14,\quad\gamma=3,\quad
 \beta=2(\gamma^2-\delta^2)=\frac{143}{8},\quad
 \eta=(\delta^2+\gamma^2)^2=\frac{21025}{256}.
 \tag{FPR1}
\]
The full ordered quartet is
\[
 w=(\delta+i\gamma,\delta-i\gamma,-\delta+i\gamma,-\delta-i\gamma),
 \qquad \rho_a=\frac12+w_a,
 \qquad V_{ar}=w_a^r\quad(0\le r\le3).
 \tag{FPR2}
\]
Let the unit be any fixed nonzero real value $a=\bar a$.
Then the original $U=\diag(a,\bar a,\bar a,a)$ is $aI$;
its cancellation in the exact conjugation below does not assign
its magnitude a new value. With $z=1/u$, PCL7--9 gives
\[
 H(u)=e^{\mathfrak a/u}G_u\mathcal R(z)V^{-1}U,
 \qquad
 M_j(z)=H(u)^{-1}D^{-j}H(u)
       =V\mathcal R(z)^{-1}D^{-j}\mathcal R(z)V^{-1},
 \tag{FPR3}
\]
where $D=\diag(\zeta,\zeta^2,\zeta^3,\zeta^4)$ and
$\zeta=e^{2\pi i/5}$. The original scalar, branch powers and
Gamma ray constants commute with $D$ and cancel in this conjugation.
They remain in $H$ itself and in any metric in those coordinates.

Put $Q_0(t)=t_1t_4-t_2t_3$ and
\[
 t(\xi)=(e^{\rho_1\xi},e^{\rho_2\xi},
                 e^{\rho_3\xi},e^{\rho_4\xi})^{\mathsf T},\quad
 f_j(\xi,z)=Q_0(M_j(z)t(\xi)),\quad
 d_j(z)=f_j(0,z),\quad E(\xi,z)=\prod_{j=1}^4 f_j(\xi,z).
 \tag{FPR4}
\]
These are exactly OFM1's original symbol factors. Here $z$ denotes
inverse period and $\xi$ denotes the symbol variable, so the two
derivatives have different mathematical meanings.
This parameter choice is a member of the geometric coefficient
family. It does not assert that the four $\rho_a$ are Riemann-zeta
zeros or that the arithmetic unit has the real phase used here.

\section{Full coefficient bound and both infinite tails}
The literal PCL2 coefficient, with $1\le a\le4$ and $0\le r\le3$, is
\[
 R_{ar,n}=
 \sum_{\substack{p,h\ge0\\2p+4h=5n+r+1-a}}
 \frac{\beta^p\eta^h}{3^p p!h!}
 (-1)^L5^L(a/5)_L,\qquad L=p+h-n.
 \tag{FPR5}
\]
For a summand let $m=p+h$ and $s=r+1-a\in[-3,3]$.
As $r+3p+h+1=a+5L$ is positive and $a\le4$, one has $L\ge0$.
Thus $m\ge n$. Because $0<a/5<1$, each factor in the rising
factorial gives $(a/5)_L\le L!$. Moreover
\[
 \frac{n!L!}{p!h!}
 =\frac{\binom{m}{p}}{\binom{m}{n}}\le2^m.
 \tag{FPR6}
\]
Choose $K=11/2$. Direct rational inequalities give
$K\ge1$, $K^2\ge5\beta/3$ and $K^4\ge5\eta$.
Consequently the absolute value of each summand in FPR5 is at most
\[
 \frac{5^{-n}}{n!}(5\beta/3)^p(5\eta)^h2^m
 \le\frac{K^{5n+3}}{5^n n!}\,4\,6^n.
 \tag{FPR7}
\]
Indeed $2p+4h=5n+s\le5n+3$, and
$m\le(5n+3)/2$, whence
$2^m\le2^{3/2}(\sqrt{32})^n\le4\cdot6^n$.
For each $h$ there is at most one $p$. The number of pairs is at
most $(5n+3)/4+1\le2n+2$. With
\[
 C=4K^3=\frac{1331}{2},\qquad
 T=\frac{6K^5}{5}=\frac{483153}{80},
 \qquad |R_{ar,n}|\le C(2n+2)\frac{T^n}{n!},
 \tag{FPR8}
\]
the Weierstrass convergence criterion proves that every original
$\mathcal R_{ar}$ is entire. In particular this proof supplies
convergence at the finite candidate without a large-period
assumption and without using a finite truncation as its justification.

For $|z|\le\varrho$, let $w=T\varrho$ and $N=400$. The omitted
series and its derivative are bounded entrywise by
\[
 \begin{split}
 B_0&=\frac{C(2N+4)w^{N+1}}{(N+1)!}
            \frac1{1-2w/(N+2)},\\
 B_1&=\frac{C(2N+4)(N+1)T w^N}{(N+1)!}
            \frac1{1-3w/(N+1)}.
 \end{split}\tag{FPR9}
\]
To prove the first bound, the successive ratio for
$(2n+2)w^n/n!$ is
$w(n+2)/(n+1)^2\le2w/(N+2)$ for $n\ge N+1$.
For the derivative summands $(2n+2)T w^{n-1}/(n-1)!$ it is
$w(n+2)/(n(n+1))\le2w/(N+1)<3w/(N+1)$.
The script verifies that both denominators are positive. Summing
the resulting geometric bounds proves FPR9, including every
omitted coefficient. Enclosing each complex tail in a rectangle
with real and imaginary radii $B_i$ is conservative and valid.

The coefficient loop enumerates every permissible $h$ from zero
through $\lfloor(5n+s)/4\rfloor$, rejects odd $5n+s-4h$, and
sets $p=(5n+s-4h)/2$. Its stored rising-factorial term is exactly
\[
 (-5)^L(a/5)_L=(-1)^L\prod_{\ell=0}^{L-1}(a+5\ell).
 \tag{FPR10}
\]
The arrays cover every such $p,h,L$ through $n=400$.
Simultaneous Horner evaluation uses the identity
$(zP+R)'=P+zP'$, hence encloses both the degree-400 series and
its derivative before adding FPR9's complete tails.

\section{The finite disk and a separate local inverse proof}
The following decimal strings designate exact rational numbers:
\[
 \begin{split}
 c_R={}&0.011136978231486808477778740306233906621055400363647396688725562781,\\
 c_I={}&0.0078780133477760021239711966284294386302928222474756316845213852718.
 \end{split}\tag{FPR11}
\]
Put $c=c_R+ic_I$, $r=10^{-14}$, and define
\[
 \mathcal D=\{z:|z-c|\le r\},\qquad
 \mathcal B=\{z:|\Re(z-c)|\le r,\ |\Im(z-c)|\le r\}.
 \tag{FPR12}
\]
The arithmetic includes the tiny rounding enclosures of $c,r$;
thus its evaluated box contains this exact $\mathcal B$, which
contains $\mathcal D$. Its $\varrho$ encloses the modulus of
every point in that box. The resulting explicit tail bounds are
$B_0<6.44\cdot10^{-98}$ and $B_1<2.92\cdot10^{-93}$.

PCL5 already supplies the global identity $\det\mathcal R=1$.
There is also a separate finite local certificate, so local
invertibility need not rest on an unexamined use of that identity.
Let $A$ be the exact dyadic midpoint matrix of the ball inverse
of the enclosed $\mathcal R(c)$. The independent checker forms
$I-A\mathcal R(\mathcal B)$ with the full tails included.
The four certified upper bounds for the row sums of entry
moduli are respectively less than
\[
 1.73\cdot10^{-9},\quad 7.75\cdot10^{-10},\quad
 1.74\cdot10^{-10},\quad 6.16\cdot10^{-11}.
 \tag{FPR13}
\]
For every exact $z\in\mathcal B$ it follows that
$\|I-A\mathcal R(z)\|_\infty<1$. Therefore
$A\mathcal R(z)$ has inverse
$\sum_{n\ge0}(I-A\mathcal R(z))^n$ by the geometric-series
identity. Both factors are square, so $\mathcal R(z)$ is
invertible. Its inverse is holomorphic near the closed box by
the adjugate formula and compactness. This proves the analytic
domain needed for the zero count.

\section{Exact derivative maps and the zero count}
For $R=\mathcal R(z)$, differentiation of FPR3 gives
\[
 M_j'(z)=V\bigl(R^{-1}D^{-j}R'
          -R^{-1}R'R^{-1}D^{-j}R\bigr)V^{-1}.
 \tag{FPR14}
\]
The order of every matrix factor in the script matches this
identity. Define $y=M_j\mathbf1$, $p=M_j'\mathbf1$ and
$x=M_j\rho$. Differentiating FPR4 yields exactly
\[
 \begin{split}
 d_j&=y_1y_4-y_2y_3,\\
 d_j'&=p_1y_4+y_1p_4-p_2y_3-y_2p_3,\\
 (\partial_\xi f_j)(0,z)&=x_1y_4+y_1x_4-x_2y_3-y_2x_3.
 \end{split}\tag{FPR15}
\]
In particular $\rho$, rather than the centered $w$, occurs in
the symbol derivative. All four labels and $D^{-j}$ signs are
the exact NC1/NC5 and OFM1 conventions.

The root script and independent replay give the strict bounds
\[
 |d_1(c)|<5\cdot10^{-63},\qquad |d_1'(c)|>434,
 \qquad \sup_{z\in\mathcal B}|d_1'(z)-d_1'(c)|<8.
 \tag{FPR16}
\]
The last enclosure subtracts the centre derivative ball from
the entire-box derivative ball, so contains every exact
derivative difference. For $|z-c|=r$, integration on the segment
from $c$ to $z$ gives
\[
 |d_1(z)-d_1'(c)(z-c)|
 \le |d_1(c)|+r\sup_{\mathcal B}|d_1'-d_1'(c)|
 <5\cdot10^{-63}+8\cdot10^{-14}
 <434\cdot10^{-14}<|d_1'(c)(z-c)|.
 \tag{FPR17}
\]
By Rouch\'e's theorem, $d_1$ and the nonzero linear polynomial
$d_1'(c)(z-c)$ have the same number of zeros, counted with
multiplicity, in $|z-c|<r$. That number is one. Hence there is
exactly one zero $z_*$ in the disk, and it is simple. The proof
uses the full holomorphic function, not just its truncation.
The same bounds also prove uniqueness on the entire rectangle.
For $z_1,z_2\in\mathcal B$, integrate along their segment to get
$|d_1(z_1)-d_1(z_2)-d_1'(c)(z_1-z_2)|\le8|z_1-z_2|$.
Thus $|d_1(z_1)-d_1(z_2)|>426|z_1-z_2|$ for distinct points,
so $d_1$ is injective there.
Since $0\notin\mathcal B$, the finite original period
$u_*=1/z_*$ exists. For orientation only its certified enclosure
is centred near $59.8455390191-42.333202544i$; FPR11--12, not
these displayed rounded coordinates, specify the exact zero.

\section{Original symbol order and the five distinct quadrics}
The complete-box enclosures additionally give
\[
 \Im\partial_\xi f_1(0,z)>3.68,\qquad
 d_2(z)\ne0,\quad d_3(z)\ne0,\quad d_4(z)\ne0
 \quad(z\in\mathcal B).
 \tag{FPR18}
\]
For example the recorded value rectangles for the last three
factors have real parts in
$[-0.23174,-0.22826]$, $[-2.42469,-2.41531]$ and
$[-3.35312,-3.34688]$, respectively; none contains zero.
At the exact $z_*$, the first factor has a simple zero in $\xi$
and the other three have order zero. Thus
\[
 \ord_{\xi=0}E(\xi,z_*)=1,\qquad
 \mu_1(z_*)=\partial_\xi f_1(0,z_*)
                    d_2(z_*)d_3(z_*)d_4(z_*)\ne0.
 \tag{FPR19}
\]
The original script encloses this complete product with modulus
strictly greater than $13.80$. Distinctly, simplicity in the
inverse-period direction gives
\[
 \frac{d}{dz}E(0,z_*)=d_1'(z_*)d_2(z_*)d_3(z_*)d_4(z_*)\ne0.
 \tag{FPR20}
\]
These two assertions concern different derivatives, and both
are needed when transporting the order jump.

Here is the entire orbit-admissibility argument. In the literal
Fourier coordinates put
\[
 Q_u(X)=Q_0(H(u)^{-1}X),\quad
 (gQ)(X)=Q(D^{-1}X),\quad Q_{u,j}=g^jQ_u\quad(0\le j\le4).
 \tag{FPR21}
\]
Because $g^5=1$, its action on association classes of nonzero
quadratic polynomials has orbit size either one or five.
More explicitly, if $Q_{u,j}=bQ_{u,k}$ with $j\ne k$ and
$b\ne0$, then $g^{j-k}[Q_u]=[Q_u]$. A nonzero residue modulo
five generates the group, so $g[Q_u]=[Q_u]$. In particular
$Q_{u,1}=b_1Q_u$ for some nonzero $b_1$.
Pulling back along $X=H(u)t(\xi)$ would give
\[
 f_1(\xi,z)=b_1Q_0(t(\xi))
 =b_1\bigl(e^{(\rho_1+\rho_4)\xi}
             -e^{(\rho_2+\rho_3)\xi}\bigr)=0,
 \tag{FPR22}
\]
since both exponent sums equal one. This contradicts FPR18.
Thus every pair of the five quadrics is nonassociate at $u_*$,
and indeed at every period reciprocal to a point of the box.
They are prime: $Q_0$ is an irreducible quadratic over $\C$
(its symmetric matrix has rank four, whereas a product of two
linear forms has rank at most two), and the invertible linear
changes $H^{-1}D^{-j}$ preserve irreducibility. The local
inverse proof establishes that these changes exist. In the
original diagonal matrix the ray constant is
$g_a=5^{a/5-1}e^{\pi ia/5}\Gamma(a/5)$, as specified in PCL3.
Its Gamma value is strictly positive by its positive-real
integral because $a/5>0$. Thus $g_a\ne0$, and both
$e^{\mathfrak a/u}$ and $u^{a/5}=\exp((a/5)\log u)$ are
nonzero on the retained logarithm branch for $u\ne0$.
The original $V$ is invertible because the four ordered
$w_a$ are distinct, and $U=aI$ is invertible because $a\ne0$.
This proves the five-prime-orbit hypothesis directly rather
than importing a sufficient large-period cutoff.

\section{Independent check of the complete weighted inverse constants}
The continuation FPZ15--24 was independently read against WCF1--6.
Here is the derivation of every degree and exterior rank, retaining
the original Gamma weights. For fixed $s\in\{1,a_{\rm amp}\}$,
$a_{\rm amp}=4\ell+1\ge9$, let
\[
 \rho_n=\sqrt{\frac{n!}{(s/2)_n}},\quad
 c_0=a_{\rm amp}/2,\quad c_1=c_0-4,\quad
 g(t,z)=e^{-4\omega(t)}E(\omega(t),z),\quad
 \omega(t)=-i\arctan t.
 \tag{FPR23}
\]
WCF2--6 gives the fixed old $v=0$ map in its actual
orthonormal source and target frames:
\[
 B_{D,s}(z)_{ij}=\begin{cases}
  g_{j-i}(z)\rho_j/\rho_i,&i\le j,\\0,&i>j,
 \end{cases}\qquad0\le i,j\le D.
 \tag{FPR24}
\]
The common original mass is $M_s=(2\pi)^{s/2}$; its two
orthonormal-frame factors cancel in this matrix, as WCF6 proves.
Neither source nor target norm or centre has changed.
At $z_*$, $g_0=0$ and $g_1=-i\mu_1\ne0$.
The other two coefficients used for $D=3$ are
$g_2=4\mu_1-\mu_2/2$ and
$g_3=i(\mu_3/6-2\mu_2+25\mu_1/3)$.
They follow by substituting
$\omega=-it+it^3/3+O(t^5)$ into FPR23, and the separate checker
verifies this formal composition exactly.

For $D\ge1$, the first column and last row of $B_*=B_{D,s}(z_*)$
are zero. The block in rows $0,\ldots,D-1$ and columns $1,\ldots,D$
is upper triangular with diagonal $g_1\rho_j/\rho_{j-1}$.
It is invertible, with determinant $g_1^D\rho_D$, since $\rho_0=1$.
Permuting the zero column to the end identifies $B_*$ with this
block plus a zero row and column by unitary coordinate permutations.
Hence its nonzero singular values are precisely the $D$ singular
values of that block, and their product is
$|\mu_1|^D\rho_D$. In particular the rank is exactly $D$, the
kernel is the constant source polynomial and the image is
$\mathcal P_{D-1}$ in the original target. For $D=0$ the matrix
is zero, with empty product one.

Set $n=D+1$ and $\alpha=\mu_0'(z_*)\ne0$. The final certificate
additionally verifies $|\alpha|>1524$.
Triangularity gives $\det B_{D,s}(z)=\mu_0(z)^n$.
Thus for $z\ne z_*$ sufficiently close, with singular values
$\sigma_1(z)\ge\cdots\ge\sigma_n(z)>0$, one has
\[
 \|\bigwedge^r B_{D,s}(z)^{-1}\|
 =\frac{\sigma_1(z)\cdots\sigma_{n-r}(z)}
        {|\mu_0(z)|^n}
 =\frac{\|\bigwedge^{n-r}B_{D,s}(z)\|}{|\mu_0(z)|^n}.
 \tag{FPR25}
\]
The numerator has a positive limit for every $1\le r\le n$,
because $n-r\le D=\operatorname{rank}B_*$. The convention for
the zeroth exterior norm is one. Since
$\mu_0(z)=\alpha(z-z_*)+O((z-z_*)^2)$, this proves
\[
 \lim_{z\to z_*}|z-z_*|^{D+1}
 \|\bigwedge^rB_{D,s}(z)^{-1}\|
 =\frac{\|\bigwedge^{D+1-r}B_*\|}{|\alpha|^{D+1}}>0.
 \tag{FPR26}
\]
Exactly $D$ forward singular values tend to positive limits.
Their product tends to $\rho_D|\mu_1|^D$, while the product
of all $D+1$ equals $|\mu_0|^{D+1}$. The last forward singular
value therefore vanishes to order $D+1$. Consequently exactly
one inverse singular value has pole $D+1$, and the other $D$
inverse singular values tend to positive finite values. The
operator-norm constant is
$\rho_D|\mu_1|^D/|\alpha|^{D+1}$.

For $D=3$ the nonzero block is precisely
\[
 C_s=\begin{pmatrix}
 g_1\rho_1&g_2\rho_2&g_3\rho_3\\
 0&g_1\rho_2/\rho_1&g_2\rho_3/\rho_1\\
 0&0&g_1\rho_3/\rho_2
 \end{pmatrix}.
 \tag{FPR27}
\]
If $\lambda_1\ge\lambda_2\ge\lambda_3>0$ are the eigenvalues
of $C_s^*C_s$, then FPR26 in exterior ranks $1,2,3,4$ gives
the four constants
$|\alpha|^{-4}(\rho_3|\mu_1|^3,
\sqrt{\lambda_1\lambda_2},\sqrt{\lambda_1},1)$.
All original moments in these entries are the finite exact
derivatives of the original four factors as in OFM5 and FPZ17;
none is a free independent replacement parameter.

Finally $z-z_*=-(u-u_*)/(uu_*)$ gives
\[
 \lim_{u\to u_*}|u-u_*|^{D+1}
 \|\bigwedge^r B_{D,s}(u^{-1})^{-1}\|
 =|u_*|^{2(D+1)}
   \frac{\|\bigwedge^{D+1-r}B_*\|}{|\alpha|^{D+1}}.
 \tag{FPR28}
\]
At the zero the actual order is $v=1$. The order-adjusted map
has domain $\mathcal P_{D+1}/\mathcal P_0$, target
$\mathcal P_D$, and nonzero diagonal
$(-i\mu_1)\rho_{j+1}/\rho_j$ by WCF5--6. It is therefore
invertible. This exact change of domain is retained alongside
the fixed-space singularity. No contradiction between these
two maps is inferred.

\section{The general finite-zero existence argument}
FPZ25--26 is also valid for every original
$0<\delta<1/2$, $\gamma>2$ on the nonzero real-unit phase.
Here the parameters are retained rather than set to FPR1.
Choose $K=\max\{1,(5\beta/3)^{1/2},(5\eta)^{1/4}\}$,
$C_R=4K^3$, $T_R=6K^5/5$. Summing FPR8 gives
\[
 |\mathcal R_{ar}(z)|\le
 H(|z|):=2C_R(1+T_R|z|)e^{T_R|z|}.
 \tag{FPR29}
\]
The sum is exact because
$\sum_{n\ge0}(2n+2)x^n/n!=2(1+x)e^x$.
For this global argument use the proved original PCL5 identity
$\mathcal R^{-1}=\operatorname{adj}\mathcal R$.
Each adjugate entry is a sum of six triple products, so
$\|\operatorname{adj}\mathcal R\|_\infty\le24H^3$
and $\|\mathcal R\|_\infty\le4H$.
Let $v_V=\|V\|_\infty$ and $v_I=\|V^{-1}\|_\infty$.
Then $\|M_j\|_\infty\le96v_Vv_IH^4$, and
$|Q_0(M_j\mathbf1)|\le2\|M_j\mathbf1\|_\infty^2$ gives
\[
 |d_j(z)|\le A(1+|z|)^8e^{B|z|},\quad
 A=2(96v_Vv_I)^2(2C_R)^8\max(1,T_R)^8,\quad B=8T_R.
 \tag{FPR30}
\]
All constants thus have finite expressions in the unchanged original
parameters. PCL22 and the parity proof following FUT7 give that $d_j$
is even. PZX/OPB15 prove its nonzero second coefficient.
FUT24--26 prove its nonzero fourth coefficient: their scalar-restored
$d_j/(4i\delta\gamma)$ has fourth coefficient whose stated imaginary
part is a nonzero scalar times $F_1(t)$ or $F_2(t)$, with
$t=\beta^2/\eta\in(0,4)$ and $h=(\sqrt5-1)/2\in(0,2/3)$.
The exact Bernstein coefficient lists there are
\[
 \begin{split}
 &(25515h,15120+21735h,16464+15939h,14848+11199h),\\
 &(25515,21735-15120h,15939-16464h,11199-14848h).
 \end{split}\tag{FPR31}
\]
Each entry is positive. For the last three entries of the second
list the lower bounds using $h<2/3$ are $11655$, $4963$ and
$3901/3$. The degree-three Bernstein basis on $[0,4]$ is
nonnegative and sums to one. Hence both polynomials are positive.
The fourth coefficients of the other two factors follow by the
proved original conjugation. The nonzero scalar $4i\delta\gamma$
is restored throughout this implication; the functions in FPR30
are the original $d_j$, not the scalar-divided presentation.

Suppose one of these original $d_j$ had no nonzero zero.
Its exact order two at zero makes $F=d_j/z^2$ entire and
zero-free. The entire $F'/F$ has a primitive on $\C$;
choosing its additive constant at zero gives an entire $h_0(z)$
such that $F=e^{h_0}$. The radial bound FPR30 and compactness
near zero yield
$\Re h_0(z)\le A_0+B|z|+8\log(1+|z|)$ for a finite $A_0$.
Write $h_0(z)=\sum_{n\ge0}b_nz^n$ and
$M_R=A_0+BR+8\log(1+R)$.
Fourier integration of the real part on $|z|=R$ gives
$b_nR^n/2=(2\pi)^{-1}\int_0^{2\pi}
\Re h_0(Re^{i\theta})e^{-in\theta}\,d\theta$ for $n\ge1$.
Subtract $M_R$ inside this nonconstant Fourier coefficient and
use the nonnegative integrand $M_R-\Re h_0$. Its integral is
$2\pi(M_R-\Re b_0)$ by the mean value identity. Consequently
\[
 |b_n|\le\frac{2(M_R-\Re b_0)}{R^n}.
 \tag{FPR32}
\]
As $R\to\infty$, this forces $b_n=0$ for $n\ge2$.
Thus $F=e^{b_0+b_1z}$. Since $d_j$ is even, so is $F$;
therefore $e^{2b_1z}=1$ for all $z$, and differentiation at
zero gives $b_1=0$. This makes $d_j=e^{b_0}z^2$, contradicting
its nonzero fourth coefficient. Hence each factor has a
nonzero finite complex zero. This argument proves existence
for every original quartet on the stated phase, but supplies
neither orbit admissibility nor the symbol order at every such
zero. Those stronger conclusions for FPR11--12 follow from
the separate enclosing calculation and FPR18--22.

\section{Review conclusion and limits of the assertion}
No mathematical error was found in the coefficient loop, the
two majorants, the original-coordinate conjugation, the two
derivative formulas or the complex Rouch\'e argument at radius
$10^{-14}$. The earlier radius $10^{-12}$ failed the interval
test; that unsuccessful test was not a proof of absence and
was not used as a positive certificate.

The certified result is a finite geometric period with a genuine
original-symbol order change from zero to one and five distinct
prime quadrics. The nonzero real unit is retained through its
exact cancellation in FPR3. The result supplies neither an
arithmetic zeta-zero quartet nor an assertion about the actual
arithmetic unit's phase. It also does not prove the separate
coordinate-section coprimality conclusions of NC19--22 or
membership in a previously imposed numerical $|u|\ge R_*$
subdomain. Those are different assertions and were not checked
by this certificate. The proved five-orbit condition itself
holds on the displayed finite box, independently of that cutoff.

\section{Exact reading and replay record}
The sources actually read for this mathematical review were:
\begin{itemize}
\item \texttt{PCL\_COMPLETE.tex}: PCL1--11 and RMR1--3;
source SHA256
\texttt{231acb7cd4902b6d3477f853d1d8dea814c71b94cff2a24e7fdee5481652ef54}.
\item \texttt{ORIGINAL\_CONDUCTOR\_COPRIMALITY.tex}: complete
file, with NC1--9 used for the original objects and NC19--22
read to delimit what the certificate does not establish;
source SHA256
\texttt{71baa0ba21e7d9de50223eeebfb12532442f270c2552cc1eea9d136629d945ab}.
\item \texttt{FABLE\_TO\_ORIGINAL\_CONDUCTOR.tex}: FC13--17
and OFM1--11 read in full;
source SHA256
\texttt{8739afeab8d76ebfa4fc5d15d8e8486d8da923beff8675a212cbbda61e00cc18}.
\item \texttt{WEIGHTED\_CONDUCTOR\_FORWARD.tex}: WCF1--12
read, with WCF1--6 used for the full weighted inverse verification;
source SHA256
\texttt{d86dcb3daf13c9b2d5a0aaa4d72cce35a1aff48066ad69c7d9846b9d097dc354}.
\item \texttt{FINITE\_PERIOD\_ZERO\_BODY.tex}: complete FPZ1--26,
including the corrected rectangle injectivity proof and the explicit
growth constant in FPZ25; reviewed SHA256
\texttt{7972cb199f16c7b8cf02332889b48bf63c078f00e2b760c66cd5bb589b650207}.
\item \texttt{FINITE\_PERIOD\_UNIT\_TEST.tex}: FUT1--7 and
the following parity proof, plus FUT20--26; FUT24--26 is used
for the full coefficient nonvanishing in FPR31; reviewed SHA256
\texttt{1d0ff101f47caba1b5f8ed741d5f38f26ee85a898607effb7b810869eb4853dd}.
\item The full root certificate script, audited and replayed
successfully with python-flint 0.9.0, 768-bit working precision
and one arithmetic thread; initial reviewed SHA256
\texttt{685598fec03bf4018c0a17ffa3a261ab6da379393dd86b44e30d0366df55a877}.
The final additional output of $\alpha$, $\mu_0,\ldots,\mu_3$
and $g_0,\ldots,g_3$ was separately read and replayed at SHA256
\texttt{041728c2ba636914c9c6dbbc294da396b7b1233b9967fb0b4af9cdb3baa8b2b8}.
The binomial factor-derivative and product rules match OFM5
and FPZ17; substituting $f_{1,0}=0$ uses the proved exact root.
\item The separate checker
\texttt{check\_finite\_period\_certificate\_review.py} and its
\texttt{FINITE\_PERIOD\_CERTIFICATE\_REVIEW\_CHECK.json}
record the extra Neumann bounds and independent rounded
Rouch\'e inequalities. They do not claim to be a second
implementation of all Arb arithmetic.
\end{itemize}
PCL, NC and OFM are credited here as prior programme calculations;
the current review does not claim authorship of them. Rouch\'e's
zero-count theorem is the classical human theorem used at FPR17.
No new human paper was claimed read in this bounded review.
\end{document}
